% GENERATED FILE. Edit canonical lessons, then run npm run generate:books.
% canonical-source-hash: fnv1a64:a3f53bf262b9268a
% canonical-lessons: FR-C302-incendie, FR-C302-extincteur, FR-C302-explosion, FR-C302-exploser, FR-C302-etincelle

\chapter{Au feu !}
\label{ch:fr-au-feu}
\hlchaptermodality{302}

\begin{chapteropening}
\textbf{By the end of this chapter:} \emph{I can describe a fire and a blast: the blaze, the extinguisher, an explosion, blowing up and a spark.}

Complete the last lesson of chapter 302: I can describe a fire and a blast: the blaze, the extinguisher, an explosion, blowing up and a spark.
\end{chapteropening}

\section[l'incendie]{l'incendie — the fire, the blaze}

\label{lesson:FR-C302-incendie}



\begin{quote}
Before the new one: say the French for the emergency exit, then the French for to evacuate.
\end{quote}

\subsection*{You'll want to know: l'incendie}
\textbf{l'incendie} — \textquotedblleft{}the fire, the blaze\textquotedblright{}.

\begin{quote}
Un incendie a brûlé trois maisons.
\end{quote}

\begin{quote}
\emph{A fire burnt down three houses.}
\end{quote}

A fire out of control is \emph{un incendie}; a fire in a fireplace is \emph{un feu}.

\begin{cousinweb}
\emph{incendie} is a learned borrowing from Latin \emph{incendium} \textquotedblleft{}a burning, a fire\textquotedblright{}, from \emph{incendere} \textquotedblleft{}to set on fire\textquotedblright{} (Lewis \& Short).
\end{cousinweb}

\begin{grammarlens}[title={What you've built}]
One more word for telling, explaining and describing at B1.
\end{grammarlens}

\subsection*{Your turn}
\begin{itemize}
\raggedright
  \item \emph{Say it:} \emph{l'incendie}
  \item \emph{Say it:} \emph{l'incendie}, once more
  \item \emph{Say it:} a sentence of your own with \emph{l'incendie}
  \item \emph{Recall:} say the French for the emergency exit, then the French for to evacuate, then say \emph{l'incendie} again
\end{itemize}

\begin{tcolorbox}[breakable,colback=teal!4,colframe=teal!35!black,title={Before you move on}]
What does l'incendie mean? (\textquotedblleft{}The fire, the blaze\textquotedblright{}.) Say it once more.
\end{tcolorbox}
\section[l'extincteur]{l'extincteur — the fire extinguisher}

\label{lesson:FR-C302-extincteur}



\begin{quote}
Before the new one: say the French for to evacuate, then the French for the fire.
\end{quote}

\subsection*{You'll want to know: l'extincteur}
\textbf{l'extincteur} — \textquotedblleft{}the fire extinguisher\textquotedblright{}.

\begin{quote}
Le cuisinier a utilisé l'extincteur.
\end{quote}

\begin{quote}
\emph{The cook used the fire extinguisher.}
\end{quote}

Masculine: \emph{un extincteur}.

\begin{cousinweb}
\emph{extincteur} is a learned borrowing from Latin \emph{exstinctor} \textquotedblleft{}one who puts out\textquotedblright{} (EtymDB), from \emph{exstinguere} \textquotedblleft{}to put out\textquotedblright{}.
\end{cousinweb}

\begin{grammarlens}[title={What you've built}]
One more word for telling, explaining and describing at B1.
\end{grammarlens}

\subsection*{Your turn}
\begin{itemize}
\raggedright
  \item \emph{Say it:} \emph{l'extincteur}
  \item \emph{Say it:} \emph{l'extincteur}, once more
  \item \emph{Say it:} a sentence of your own with \emph{l'extincteur}
  \item \emph{Recall:} say the French for to evacuate, then the French for the fire, then say \emph{l'extincteur} again
\end{itemize}

\begin{tcolorbox}[breakable,colback=teal!4,colframe=teal!35!black,title={Before you move on}]
What does l'extincteur mean? (\textquotedblleft{}The fire extinguisher\textquotedblright{}.) Say it once more.
\end{tcolorbox}
\section[l'explosion]{l'explosion — the explosion, the blast}

\label{lesson:FR-C302-explosion}



\begin{quote}
Before the new one: say the French for the fire, then the French for the fire extinguisher.
\end{quote}

\subsection*{You'll want to know: l'explosion}
\textbf{l'explosion} — \textquotedblleft{}the explosion, the blast\textquotedblright{}.

\begin{quote}
On a entendu l'explosion à deux kilomètres.
\end{quote}

\begin{quote}
\emph{People heard the blast two kilometres away.}
\end{quote}

Figurative: \emph{une explosion de joie} — an outburst of joy.

\begin{cousinweb}
\emph{explosion} is a learned borrowing from Latin \emph{explōsiō} (EtymDB), from \emph{explōdere} \textquotedblleft{}to drive off by clapping\textquotedblright{}, built on \emph{plaudere} \textquotedblleft{}to clap\textquotedblright{}.
\end{cousinweb}

\begin{grammarlens}[title={What you've built}]
One more word for telling, explaining and describing at B1.
\end{grammarlens}

\subsection*{Your turn}
\begin{itemize}
\raggedright
  \item \emph{Say it:} \emph{l'explosion}
  \item \emph{Say it:} \emph{l'explosion}, once more
  \item \emph{Say it:} a sentence of your own with \emph{l'explosion}
  \item \emph{Recall:} say the French for the fire, then the French for the fire extinguisher, then say \emph{l'explosion} again
\end{itemize}

\begin{tcolorbox}[breakable,colback=teal!4,colframe=teal!35!black,title={Before you move on}]
What does l'explosion mean? (\textquotedblleft{}The explosion, the blast\textquotedblright{}.) Say it once more.
\end{tcolorbox}
\section[exploser]{exploser — to explode, to blow up}

\label{lesson:FR-C302-exploser}



\begin{quote}
Before the new one: say the French for the fire extinguisher, then the French for the explosion.
\end{quote}

\subsection*{You'll want to know: exploser}
\textbf{exploser} — \textquotedblleft{}to explode, to blow up\textquotedblright{}.

\begin{quote}
Une bouteille de gaz a explosé.
\end{quote}

\begin{quote}
\emph{A gas bottle exploded.}
\end{quote}

Informal: \emph{exploser de rire} — to burst out laughing.

\begin{cousinweb}
\emph{exploser} is built on \emph{explosion}, a learned borrowing from Latin \emph{explōsiō}, from \emph{explōdere} (Lewis \& Short).
\end{cousinweb}

\begin{grammarlens}[title={What you've built}]
One more word for telling, explaining and describing at B1.
\end{grammarlens}

\subsection*{Your turn}
\begin{itemize}
\raggedright
  \item \emph{Say it:} \emph{exploser}
  \item \emph{Say it:} \emph{exploser}, once more
  \item \emph{Say it:} a sentence of your own with \emph{exploser}
  \item \emph{Recall:} say the French for the fire extinguisher, then the French for the explosion, then say \emph{exploser} again
\end{itemize}

\begin{tcolorbox}[breakable,colback=teal!4,colframe=teal!35!black,title={Before you move on}]
What does exploser mean? (\textquotedblleft{}To explode, to blow up\textquotedblright{}.) Say it once more.
\end{tcolorbox}
\section[l'étincelle]{l'étincelle — the spark}

\label{lesson:FR-C302-etincelle}



\begin{quote}
Before the new one: say the French for the explosion, then the French for to explode.
\end{quote}

\subsection*{You'll want to know: l'étincelle}
\textbf{l'étincelle} — \textquotedblleft{}the spark\textquotedblright{}.

\begin{quote}
Une petite étincelle a suffi.
\end{quote}

\begin{quote}
\emph{One small spark was enough.}
\end{quote}

Figurative too: \emph{l'étincelle d'une idée}.

\begin{cousinweb}
\emph{étincelle} is inherited from Vulgar Latin \emph{stincilla} \textquotedblleft{}spark\textquotedblright{} (EtymDB), an altered form of Latin \emph{scintilla}.
\end{cousinweb}

\begin{grammarlens}[title={What you've built}]
One more word for telling, explaining and describing at B1.
\end{grammarlens}

\subsection*{Your turn}
\begin{itemize}
\raggedright
  \item \emph{Say it:} \emph{l'étincelle}
  \item \emph{Say it:} \emph{l'étincelle}, once more
  \item \emph{Say it:} a sentence of your own with \emph{l'étincelle}
  \item \emph{Recall:} say the French for the explosion, then the French for to explode, then say \emph{l'étincelle} again
\end{itemize}

\begin{tcolorbox}[breakable,colback=teal!4,colframe=teal!35!black,title={Before you move on}]
What does l'étincelle mean? (\textquotedblleft{}The spark\textquotedblright{}.) Say it once more.
\end{tcolorbox}
